\chapter{Tutorien}

\section{Tutorium 7 – Linearität, Zeitinvarianz, BIBO-Stabilität}

% ================================================================
\subsection*{Teil A – Linearität und Zeitinvarianz}
% ================================================================

Für jedes System wird geprüft:
\begin{itemize}
  \item \textbf{Linearität} (Abschnitt~\ref{sec:zeitinvarianz-linearitaet}):
  $\mathcal{S}[a\,x_1(t) + b\,x_2(t)] \stackrel{!}{=} a\,y_1(t) + b\,y_2(t)$
  \item \textbf{Zeitinvarianz} (Abschnitt~\ref{sec:zeitinvarianz-linearitaet}):
  $\mathcal{S}[x(t - \tau_0)] \stackrel{!}{=} y(t - \tau_0)$
\end{itemize}

% ----------------------------------------------------------------
\subsubsection*{A1 – $y(t) = x(t) + a$, \quad $a \in \mathbb{R}$, $a \neq 0$}
% ----------------------------------------------------------------

\begin{beispiel}[Linearität und Zeitinvarianz: $y(t) = x(t) + a$]

\textbf{Gegeben:} $\mathcal{S}[x(t)] = x(t) + a$,\quad $a \in \mathbb{R}$, $a \neq 0$.

\textbf{Gesucht:} Ist das System linear? Ist es zeitinvariant?

\textbf{Formel} (Abschnitt~\ref{sec:zeitinvarianz-linearitaet}): Superpositionsprinzip
und Verschiebungsinvarianz.

\textbf{Rechnung – Linearität:}

Sei $y_1(t) = x_1(t) + a$ und $y_2(t) = x_2(t) + a$. Dann:

$$
\mathcal{S}[a_1 x_1(t) + a_2 x_2(t)]
= a_1 x_1(t) + a_2 x_2(t) + a
$$

$$
a_1 y_1(t) + a_2 y_2(t)
= a_1(x_1(t)+a) + a_2(x_2(t)+a)
= a_1 x_1(t) + a_2 x_2(t) + (a_1+a_2)\,a
$$

Da $(a_1+a_2)\,a \neq a$ im Allgemeinen: $\mathcal{S}[a_1 x_1 + a_2 x_2] \neq a_1 y_1 + a_2 y_2$.

\textbf{Rechnung – Zeitinvarianz:}

$$
\mathcal{S}[x(t-\tau_0)] = x(t-\tau_0) + a
\qquad
y(t-\tau_0) = x(t-\tau_0) + a
$$

Beide Seiten sind identisch.

\textbf{Lösung:}

$$
\boxed{\text{nicht linear},\quad \text{zeitinvariant (NL-Z)}}
$$

Der konstante Offset $a$ verletzt das Superpositionsprinzip (Homogenitätsbedingung: $\mathcal{S}[0] = a \neq 0$).

\end{beispiel}

% ----------------------------------------------------------------
\subsubsection*{A2 – $y(t) = t\,x(t)$}
% ----------------------------------------------------------------

\begin{beispiel}[Linearität und Zeitinvarianz: $y(t) = t\,x(t)$]

\textbf{Gegeben:} $\mathcal{S}[x(t)] = t\,x(t)$.

\textbf{Gesucht:} Ist das System linear? Ist es zeitinvariant?

\textbf{Formel} (Abschnitt~\ref{sec:zeitinvarianz-linearitaet}).

\textbf{Rechnung – Linearität:}

$$
\mathcal{S}[a_1 x_1(t) + a_2 x_2(t)]
= t\bigl(a_1 x_1(t) + a_2 x_2(t)\bigr)
= a_1\,t\,x_1(t) + a_2\,t\,x_2(t)
= a_1 y_1(t) + a_2 y_2(t)
$$

Das Superpositionsprinzip ist erfüllt.

\textbf{Rechnung – Zeitinvarianz:}

$$
\mathcal{S}[x(t-\tau_0)] = t\,x(t-\tau_0)
$$

$$
y(t-\tau_0) = (t-\tau_0)\,x(t-\tau_0)
$$

Da $t \neq t - \tau_0$ für $\tau_0 \neq 0$: $\mathcal{S}[x(t-\tau_0)] \neq y(t-\tau_0)$.

\textbf{Lösung:}

$$
\boxed{\text{linear},\quad \text{nicht zeitinvariant (L-NZ)}}
$$

Der zeitabhängige Koeffizient $t$ bewirkt, dass eine Verschiebung des Eingangs
nicht zur gleichen Verschiebung des Ausgangs führt. Das System ist \textbf{zeitvariant}.

\end{beispiel}

% ----------------------------------------------------------------
\subsubsection*{A3 – $y(t) = \dfrac{\mathrm{d}}{\mathrm{d}t}\,x(t)$}
% ----------------------------------------------------------------

\begin{beispiel}[Linearität und Zeitinvarianz: $y(t) = \dot{x}(t)$]

\textbf{Gegeben:} $\mathcal{S}[x(t)] = \dfrac{\mathrm{d}}{\mathrm{d}t}\,x(t)$.

\textbf{Gesucht:} Ist das System linear? Ist es zeitinvariant?

\textbf{Formel} (Abschnitt~\ref{sec:zeitinvarianz-linearitaet}).

\textbf{Rechnung – Linearität:}

Die Ableitung ist ein linearer Operator (Linearität der Differentiation):

$$
\mathcal{S}[a_1 x_1(t) + a_2 x_2(t)]
= \frac{\mathrm{d}}{\mathrm{d}t}\bigl(a_1 x_1(t) + a_2 x_2(t)\bigr)
= a_1\,\dot{x}_1(t) + a_2\,\dot{x}_2(t)
= a_1 y_1(t) + a_2 y_2(t)
$$

\textbf{Rechnung – Zeitinvarianz:}

$$
\mathcal{S}[x(t-\tau_0)]
= \frac{\mathrm{d}}{\mathrm{d}t}\,x(t-\tau_0)
= \dot{x}(t-\tau_0)
= y(t-\tau_0)
$$

(Kettenregel: $\frac{\mathrm{d}}{\mathrm{d}t}x(t-\tau_0) = \dot{x}(t-\tau_0)\cdot 1$)

\textbf{Lösung:}

$$
\boxed{\text{linear},\quad \text{zeitinvariant} \;\Rightarrow\; \textbf{LTI-System}}
$$

Der Differentiationsoperator ist linear und verschiebt das Signal ohne
zusätzliche Zeitabhängigkeit — das System ist ein LTI-System.

\end{beispiel}

% ================================================================
\subsection*{Teil B – BIBO-Stabilität}
% ================================================================

Kriterium (Abschnitt~\ref{sec:faltung-eigenschaften}): Ein zeitkontinuierliches LTI-System
mit Impulsantwort $h(t)$ ist BIBO-stabil genau dann, wenn

$$
\int_{-\infty}^{\infty} |h(\xi)|\,\mathrm{d}\xi < \infty\,.
$$

% ----------------------------------------------------------------
\subsubsection*{B1 – $h(t) = 2\,\varepsilon(t-2T)\,\sin\!\left(\dfrac{\pi(t-2T)}{2T}\right)$}
% ----------------------------------------------------------------

\begin{beispiel}[BIBO-Stabilität: ungedämpfte Sinusantwort]

\textbf{Gegeben:}

$$
h(t) = 2\,\varepsilon(t-2T)\,\sin\!\left(\frac{\pi(t-2T)}{2T}\right)
$$

\textbf{Gesucht:} Ist das System BIBO-stabil?

\textbf{Formel} (Abschnitt~\ref{sec:faltung-eigenschaften}):
$\displaystyle\int_{-\infty}^{\infty} |h(\xi)|\,\mathrm{d}\xi \stackrel{?}{<} \infty$

\textbf{Rechnung:}

Substitution $\tau = t - 2T$ (der Sprung verschiebt den Träger auf $\tau \geq 0$):

$$
\int_{-\infty}^{\infty} |h(t)|\,\mathrm{d}t
= \int_{0}^{\infty} 2\left|\sin\!\left(\frac{\pi\tau}{2T}\right)\right|\mathrm{d}\tau
$$

Der Integrand $2\left|\sin\!\left(\tfrac{\pi\tau}{2T}\right)\right|$ ist periodisch mit Periode
$2T$ und hat konstante Amplitude $2$ — er klingt \textbf{nicht} ab. Das Integral über $[0,\infty)$
wächst ohne Schranke:

$$
\int_{0}^{2nT} 2\left|\sin\!\left(\frac{\pi\tau}{2T}\right)\right|\mathrm{d}\tau
= n \cdot \frac{8T}{\pi} \;\xrightarrow{n\to\infty}\; \infty
$$

\textbf{Lösung:}

$$
\boxed{\int_{-\infty}^{\infty}|h(t)|\,\mathrm{d}t = \infty \;\Rightarrow\; \textbf{nicht BIBO-stabil}}
$$

Eine ungedämpfte Schwingung in der Impulsantwort führt stets zur Instabilität, da
der Energieinhalt unbegrenzt wächst.

\end{beispiel}

% ----------------------------------------------------------------
\subsubsection*{B2 – $h(t) = 4A\!\left[\operatorname{rect}\!\left(\dfrac{t}{T}\right) + e^{-(t-T)/T}\,\varepsilon(t-T)\right]$}
% ----------------------------------------------------------------

\begin{beispiel}[BIBO-Stabilität: Rechteck + abklingende Exponentialfunktion]

\textbf{Gegeben:}

$$
h(t) = 4A\!\left[\operatorname{rect}\!\left(\frac{t}{T}\right)
       + e^{-(t-T)/T}\,\varepsilon(t-T)\right],
\quad A > 0,\; T > 0
$$

\textbf{Gesucht:} Ist das System BIBO-stabil?

\textbf{Formel} (Abschnitt~\ref{sec:faltung-eigenschaften}):
$\displaystyle\int_{-\infty}^{\infty}|h(\xi)|\,\mathrm{d}\xi \stackrel{?}{<} \infty$

\textbf{Rechnung:}

Die beiden Terme haben disjunkte Träger und können separat integriert werden.

\textbf{Term 1 –} $\operatorname{rect}(t/T)$: Träger $|t| \leq T/2$, Amplitude $4A$:

$$
\int_{-\infty}^{\infty} 4A\left|\operatorname{rect}\!\left(\frac{t}{T}\right)\right|\mathrm{d}t
= 4A \int_{-T/2}^{T/2}\mathrm{d}t = 4AT
$$

\textbf{Term 2 –} $e^{-(t-T)/T}\,\varepsilon(t-T)$: Träger $t \geq T$, Substitution $u = (t-T)/T$:

$$
\int_{-\infty}^{\infty} 4A\,e^{-(t-T)/T}\,\varepsilon(t-T)\,\mathrm{d}t
= 4A\int_{T}^{\infty} e^{-(t-T)/T}\,\mathrm{d}t
= 4AT\int_{0}^{\infty} e^{-u}\,\mathrm{d}u
= 4AT
$$

\textbf{Gesamt:}

$$
\int_{-\infty}^{\infty}|h(t)|\,\mathrm{d}t = 4AT + 4AT = 8AT < \infty
$$

\textbf{Lösung:}

$$
\boxed{\int_{-\infty}^{\infty}|h(t)|\,\mathrm{d}t = 8AT < \infty \;\Rightarrow\; \textbf{BIBO-stabil}}
$$

Das Rechteck ist zeitlich begrenzt (endliche Energie), die Exponentialfunktion klingt
für $t \geq T$ monoton ab — beide Terme liefern endliche Beiträge.

\end{beispiel}
